\documentclass[openany,fontset=fandol]{ctexbook}
\usepackage{amsmath,amssymb}
\usepackage[a4paper,margin=25mm]{geometry}
\usepackage[unicode,hidelinks]{hyperref}
\setlength{\emergencystretch}{2em}
\raggedbottom
\newcommand{\PD}{\check P}
\newcommand{\cod}{\operatorname{codim}}
\newcommand{\Sing}{\operatorname{Sing}}
\begin{document}
\begin{center}
{\Large\bfseries 代数几何基础第五卷预备笔记}\par\medskip
{\LARGE\bfseries 几何正则性的例外超平面}\par\medskip
第108批：命题8.13\quad 2026年11月24日
\end{center}

\section*{底本和约定}
法稿物理45页自命题8.13标题起，续物理46页，止于推论8.14前；对读\href{https://webusers.imj-prg.fr/~leila.schneps/grothendieckcircle/EGA/EGA5trans.pdf}{英编稿}物理31页（另用32页核对后文）。下文用 $k$ 表基域，改用 $n$ 表原稿条件 $(R_k)$ 的下标。$(R_j)$ 表示余维不超过 $j$ 的点处局部环正则；几何 $(R_j)$ 要在代数闭扩张上检验。设 $X'$ 为 $X/k$ 的光滑开集，$T=X\setminus X'$；$Y_\xi$ 是由参数 $\xi\in\PD$ 确定的超平面截面，$Y'_\xi$ 是其在 $X'$ 上的开部分，$T_\xi$ 是 $T$ 在 $Y_\xi$ 上的逆像。对偶空间中的余维取闭包后计算。几何分支和奇异点均在扩域后理解。

\section*{命题8.13：经译校的陈述}
仍设 $f:X\to P$ 无分歧，且 $X$ 无孤立点；设 $X$ 几何满足 $(R_n)$。为使下述含截面正则性要求的结论成立，另要求 $X$ 在几何基变换后没有零维伴随点。\footnote{末句为译校补足，不在法稿45页可辨的增补中。原证依赖8.2；其准确前提涉及全部伴随点，不只是底空间中孤立的点。若仅字面要求无孤立点，见末页反例。几何约化且无零维分支时，不需额外验证此条件。}
令 $Z_j$ 为 $\PD$ 中不满足以下联合条件的参数集\footnote{法46页打字清楚给出条件 $R_k$，英31误作 $E_k$。此处“正则截面”指局部方程是非零因子，而不是说其零点概形光滑。}：
\[
 \varphi_\xi\ \text{在 }X_{k(\xi)}\text{ 上为正则截面},
 \qquad Y_\xi\ \text{几何满足 }(R_j).
\]
于是：
\begin{enumerate}
\item[(a)] 若 $X'$ 的所有维数至多 $n$ 的几何不可约分支均相对于 $f$ 为直纹型，则 $\cod(Z_{n-1},\PD)\ge2$。
\item[(b)] 若 $X$ 的每个维数不超过 $n+1$ 的几何不可约分支 $X_i$ 都由 $X/k$ 的光滑点组成，并且相对于 $f$ 是直纹型的，则 $\cod(Z_n,\PD)\ge2$。\footnote{法46页读作不超过 $k+1$，英31作 $k-1$。截面使维数降一，故需控制至 $n+1$ 维的分支；英32推论8.14(b)对几何正常截面要求控制至二维分支，也与 $n=1$ 时的 $n+1$ 一致。末页有检验该差别的例子。}
\end{enumerate}
这里“直纹型”沿用8.12：该分支在光滑处的切参数之像的闭包余维至少二，不要求抽象簇带有给定的直纹结构。

\newpage
\section*{证明译文}
事实上，对每个 $\xi$，$Y_\xi$ 的几何奇异集是 $\Sing(Y'_\xi)$ 与 $T_\xi$ 的并。这里限于 $\varphi_\xi$ 正则的参数；据8.2，这只舍弃余维至少二的参数。故该奇异集在 $Y_\xi$ 中的余维等于
\[
 \min\bigl\{\cod(\Sing(Y'_\xi),Y'_\xi),
                   \ \cod(T_\xi,Y_\xi)\bigr\}.
\]
只考虑 $\Sing(Y'_\xi)$ 有限的 $\xi$；这也无妨，因为这样仍只舍弃余维至少二的集合。此时 $Y'_\xi$ 的几何奇异点都是孤立的。由此及8.4可以得到结论。

\section*{证明补足：光滑开集上的奇异点}
\paragraph{两个奇异集合的并。}
在截面经过的点处，正则截面给出一个非零因子 $h\in\mathfrak m_A$。在代数闭扩域后，若 $A/(h)$ 正则，则 $A$ 正则（\href{https://stacks.math.columbia.edu/tag/00NU}{Stacks 00NU}）；也可由
\[
 \dim A/(h)=\dim A-1,\qquad
 \operatorname{edim}A\le \operatorname{edim}(A/(h))+1
\]
看出。因此，几何不光滑的环境点一旦进入截面，不会成为截面的光滑点。另一方面，在 $X'$ 上只需直接检验 $Y'_\xi$。局部环在开子概形中不变，于是得到上面的余维公式，约定空集余维为 $+\infty$。

\paragraph{为什么有限奇异集是可用的限制。}
在一个 $d$ 维几何光滑分支上，切超平面由余法空间给出，每点的参数纤维为 $\mathbb P^{r-d-1}$（$d=r$ 时为空），故非空切关联的总维数是 $r-1$，其中 $r=\dim P$。对偶参数中几何纤维至少一维的集合因而维数不超过 $r-2$：在每个这样的不可约参数分层上，关联的逆像至少比参数多一维。取有限个分支的例外集之并，舍弃后各奇异集均为零维或空。若该分支是直纹型，还可以连它全部切参数一起舍弃，仍只舍弃余维至少二的闭集。

因此在 (a) 中，剩下可能出现孤立奇异点的光滑分支维数 $d\ge n+1$；这类点在截面中的余维为 $d-1\ge n$。在 (b) 中，剩下的光滑分支维数 $d\ge n+2$；相应余维为 $d-1\ge n+1$。这解释了两个下标和维数界相差一的原因。

\newpage
\section*{证明补足：原概形的奇异闭集}
由几何 $(R_n)$，$\cod(T,X)\ge n+1$。在 (a) 中，8.4的基本截面余维估计直接给出 $\cod(T_\xi,Y_\xi)\ge n$。在 (b) 中，同一估计只能给出 $\ge n$，还需排除等号发生在余维一的参数族上。8.4给出的唯一可能来源是：$T$ 的某个几何不可约分支在 $X$ 中余维恰为 $n+1$，且其 $f$-像零维。由于 $f$ 无分歧，因而局部拟有限，有限型情形下这样的分支只能是 $T$ 的一个孤立闭点 $t$，并有 $\dim\mathcal O_{X,t}=n+1$。但这样的 $t$ 至少位于一个维数至多 $n+1$ 的几何分支，而 (b) 要求该分支上的点均在 $X'$，矛盾。于是除去余维至少二的参数，$\cod(T_\xi,Y_\xi)\ge n+1$。与光滑开集上的估计合并，分别得到几何 $(R_{n-1})$ 和 $(R_n)$。

\paragraph{补足8.2所用的正则截面前提。}
在代数闭域上，非正则截面的局部方程落入某个伴随点对应的素理想，等价于超平面包含该伴随点闭包的像。若各伴随点的闭包维数均为正，局部拟有限性保证其像的闭包维数也为正。包含其中任一像的超平面须包含它的线性张成，因而在 $\PD$ 中至少给出两个独立线性约束。伴随点有限，故这些例外参数合起来仍为余维至少二。这是此处对8.2的应用；不能把“无孤立点”直接替换成“无零维伴随点”。同一论证可在几何基变换上检验，正则截面条件沿域的忠实平坦扩张下降。

\section*{划改记录}
\begin{itemize}
\item 法45页：框内被大叉划去的旧开头曾列 $X'$ 并对切关联像加条件，后接特征零及8.11；不是可直接并入现陈述的额外假设。正文从可辨的未划去部分与手写替代读出无分歧、无孤立点及几何 $(R_n)$。
\item 法46页：两条原打字稿中的“必要性”被划改，正文按可辨的修订层只译充分性；证明后另有拟改成充要条件的大段草稿，被整段括划，不作为命题结论。
\item 法46页：用可辨的手写替代“$\Sing(Y'_\xi)$ 有限”校正旧打字层；原稿在证明附近留有边注和改写线，未确证字样不补造。英编稿31页标出的第48注也表明其编辑读法有待核。
\end{itemize}

\newpage
\section*{两个译校检验（非原文）}
\paragraph{一、维数上界不能读作 $n-1$。}
取特征零代数闭域上的光滑二次曲面
\[
 X=\{x_0x_3-x_1x_2=0\}\subset\mathbb P^3,\qquad n=1.
\]
把 $X$ 写成秩一的 $2\times2$ 矩阵的射影簇。梯度给出由可逆线性变换诱导的对偶点，故切平面参数的像是对偶空间中的二次曲面，余维为一。在点 $[1:0:0:0]$ 的切平面 $x_3=0$ 上，截面方程为 $x_1x_2=0$，两条不同直线相交。左右换基可把任意秩一矩阵送到该点，所以每个切平面截面都有此类交点，交点是曲线上的余维一点，其局部环不正则。这意味着 $Z_1$ 包含对偶二次曲面。若按英编稿读成“维数不超过 $n-1$”，本例中此假设为空，但结论不成立。法稿的 $n+1$ 则会检验这个二维分支，并排除此例。

\paragraph{二、底空间无孤立点不排除零维伴随点。}
取 $n\ge1$、$d\ge n+2$，在 $\mathbb P^{d+1}$ 中用齐次坐标 $[u:x_1:\cdots:x_d:z]$，令 $X$ 由理想
\[
 I=(zx_1,\ldots,zx_d,z^2)
\]
定义，并令 $f$ 为闭浸入。$X$ 的底空间是 $z=0$ 的 $\mathbb P^d$，没有孤立点。在 $p=[1:0:\cdots:0]$ 之外，某个 $x_i$ 可逆，从而 $z=0$，所以 $X$ 光滑。唯一例外点 $p$ 的余维为 $d>n$，因此 $X$ 几何满足 $(R_n)$；而 (a)、(b) 的低维分支假设都是空的。

但在 $u=1$ 的坐标环中，$z$ 的类非零，且被 $x_1,\ldots,x_d,z$ 全部消去。它在 $p$ 的局部环中仍非零：其消去理想恰为极大理想，局部化所倒置的元素无一在此理想中。每个经过 $p$ 的超平面局部方程都在这个极大理想中，因而消去该非零类，绝非正则截面。经过 $p$ 的超平面在 $\PD$ 中组成一个完整的余维一超平面。于是两条余维至少二的结论都不能仅凭未补足的普通字面假设推出。

\paragraph{接续位置。}
下一段从法稿物理46页下部“Combinant 8.13”及推论8.14开始，接法47页，对读英编稿32页。保留归约、深度条件以及8.15自我保留的草稿层，尤其复核推论引用8.13时是否需要另写伴随点条件。1158页合订和正式印刷第四卷均未修改。
\end{document}
